\chapter{Conclusions and Recommendations}

\section{Conclusions}

This thesis presented a comprehensive, open-source end-to-end computational
framework for the binary classification of EGFR inhibitor activity and the
prospective virtual screening of novel lead compounds for drug-resistant NSCLC.
The following principal conclusions can be drawn:

\begin{enumerate}
    \item \textbf{Data quality and curation are foundational.} Applying rigorous
          ChEMBL curation---including filtering for exact IC$_{50}$ values, pIC$_{50}$
          conversion, RDKit-based salt removal, canonicalisation, and molecular weight
          filtering---produced a high-quality, unambiguous dataset of 10,523 unique
          EGFR-bioactive compounds suitable for deep learning.

    \item \textbf{Fingerprint baselines outperformed the graph neural networks.}
          Under strict Bemis-Murcko scaffold evaluation, the Random Forest baseline
          achieved the highest test AUROC (0.8827), ahead of the MLP on ECFP4
          fingerprints (0.8547) and the five graph neural networks (0.8087--0.8454).
          Among the GNNs, GCN was strongest at default settings, while
          AttentiveFP---which natively exploits the 29-dimensional atom and
          12-dimensional bond features---required Bayesian hyperparameter optimisation
          to reach mid-field on the validation set (best validation AUROC 0.8632),
          although its test AUROC (0.7970) remained the lowest of the GNNs. These
          results show that, on a medium-sized scaffold-split EGFR dataset, well-tuned
          fingerprint models remain more accurate than graph neural networks for
          classification.

    \item \textbf{Nepali medicinal plant virtual screening is model-sensitive but
          converges on flavonoid leads.} Screening the 535-compound library with the
          two best models gave sharply different funnels---the best classifier
          (Random Forest) predicted 181 actives ($\hat{p} \geq 0.50$) but only one
          compound above $\hat{p} \geq 0.90$, while the best graph network (GCN)
          predicted 24 actives and none above $\hat{p} \geq 0.90$---exposing the
          earlier tuned-AttentiveFP screen (71 compounds above $\hat{p} \geq 0.90$) as
          over-confident. Despite this divergence, both best models enriched for the
          same drug-like chemotype: the flavones \textbf{apigenin} (\textit{Ocimum
          sanctum}; Random Forest $\hat{p} = 0.913$, QED $= 0.632$), \textbf{kaempferol},
          and \textbf{luteolin}. Luteolin, predicted active by \emph{both} best models,
          is the top cross-model consensus candidate for experimental validation.

    \item \textbf{Consensus across models is essential for trustworthy natural-product
          screening.} Because the predicted hit list depends strongly on the choice of
          classifier and the models are not calibrated on out-of-distribution natural
          products, no single model can be treated as authoritative. Prioritising
          compounds endorsed by more than one strong model---and coupling activity
          prediction with drug-likeness filtering, since the weaker models still
          ranked large, non-drug-like \textit{Terminalia chebula} tannins and
          \textit{Tinospora} triterpenoids among their highest scorers---is the key
          methodological lesson of the screening campaign. The flavone leads that survive this scrutiny, combined
          with the low cost and cultural accessibility of their source plants in Nepal,
          make them valuable starting points for affordable NSCLC drug discovery.

    \item \textbf{The pipeline is reproducible, scalable, and regionally relevant.}
          The entire framework uses exclusively open-source Python tools (PyTorch
          Geometric, RDKit, Optuna, scikit-learn, COCONUT 2.0) and public databases
          (ChEMBL, COCONUT), runs on commodity CPU hardware, and is immediately
          adaptable to other kinase targets or plant species without specialised
          hardware or licensed software. This makes it directly accessible to academic
          researchers in Nepal and other resource-limited South Asian settings.
\end{enumerate}

\section{Recommendations for Future Work}

\begin{enumerate}
    \item \textbf{Experimental validation of Nepali plant leads.} The top drug-like
          natural-product leads---particularly the flavones apigenin
          (\textit{Ocimum sanctum}), luteolin, and kaempferol, prioritised by the best
          models and cross-model consensus---should be subjected to
          \textit{in vitro} EGFR kinase inhibition assays and EGFR-mutant NSCLC cell
          line proliferation assays (e.g., PC9, HCC827) to validate the computational
          predictions empirically.

    \item \textbf{Mutant-specific modelling.} As mutation-specific bioactivity datasets
          for T790M and C797S EGFR grow in public databases, dedicated GNN models trained
          specifically on these mutant targets should be developed to directly address
          drug resistance.

    \item \textbf{3D graph representation.} Incorporating three-dimensional molecular
          conformations (e.g., through SchNet or DimeNet++ architectures) and protein
          structure-aware approaches could improve the prediction of binding affinity
          for specific EGFR mutant forms.

    \item \textbf{Pre-training on large molecular datasets.} Applying self-supervised
          graph pre-training on millions of unlabelled molecules (e.g., from ZINC or
          PubChem) before fine-tuning on EGFR data could further improve scaffold
          generalisation.

    \item \textbf{Broader Nepali plant species survey.} The current study covered
          seven ethnopharmacologically relevant species. A systematic expansion to the
          full catalogue of 700+ documented Nepali medicinal plants, using COCONUT 2.0
          and complementary databases (e.g., Dr.\ Duke's Phytochemical Database, NPACT),
          could identify additional EGFR inhibitor candidates from the rich Himalayan
          biodiversity.

    \item \textbf{Extended ADMET profiling.} The current ADMET analysis used RDKit
          physicochemical property proxies. Future work should incorporate dedicated
          ADMET prediction tools such as SwissADME, pkCSM, or deep ADMET models for
          comprehensive absorption, distribution, metabolism, excretion, and toxicity
          profiling of the top leads.

    \item \textbf{Collaboration with Nepal-based research institutions.} To translate
          these computational predictions into experimental evidence, partnerships should
          be established with Nepal-based pharmacology and oncology research institutions---such
          as the B.P.\ Koirala Institute of Health Sciences (BPKIHS) and the National
          Cancer Hospital and Research Centre (NCHRC)---for wet-lab assay, phytochemical
          isolation, and clinical-translational validation of the top Nepali medicinal
          plant lead compounds.
\end{enumerate}
